% Volume I: The Epistemology of Suspicion
% Part I. The Power of Paranoia
% Chapter 3: The First Available Story
% Spine: proof-spines/ch003-spine.md

\chapter{The First Available Story}
\label{ch:ch003}


If interpretations arrive sequentially as \(H_1,H_2,\ldots,H_n\), the first available story acquires advantages before truth has been settled. Early narratives benefit from anchoring, repetition, indexing, search ranking, organizational memory, and resource allocation. Let \(w_i(t)\) be the institutional weight assigned to hypothesis \(H_i\). A simple reinforcement process is
\[
w_i(t+1)
=
(1-\lambda)w_i(t)
+
\lambda s_i(t)
+
\rho a_i(t),
\]
where \(s_i(t)\) is evidentiary support and \(a_i(t)\) is accumulated attention. When \(\rho\) is large, attention becomes partially self-validating: uptake helps create the very standing later interpreted as confirmation.

The chapter names this persistence \textbf{narrative hysteresis}. Even when the underlying evidence changes, public judgment can remain trapped in an earlier interpretive basin because the first story has already reorganized memory, testimony, and institutional effort around itself. The task of disciplined suspicion here is not to deny narrative need, but to resist treating temporal precedence as epistemic entitlement.
